\section{The coloring reduction and the unique sixty-one-vertex residual}
\label{sec:frontend}

\subsection{The minimum-degree pivot}

Let $G$ be the least-edge color graph from \cref{lem:caps}. If $\delta(G)\geq12$, then the handshake lemma and $e(G)\leq870$ force $G$ to be $12$-regular. Its maximum degree is $12$, and it has no $K_{13}$ component because $\omega(G)\leq12$. Brooks' theorem therefore gives a proper vertex coloring of $G$ with $12$ colors \cite{Brooks1941}. One color class has at least $\lceil145/12\rceil=13$ vertices and is independent, contradicting $\alpha(G)\leq12$. Hence a minimum-degree vertex $v$ has degree
\[
d=d_G(v)\in\{0,1,\ldots,11\}.
\]
Put $A=N_G(v)$ and $U=V(G)\setminus N_G[v]$, so $|U|=144-d$.

\subsection{Packing target cliques}

Choose an inclusion-maximal family of $j$ pairwise vertex-disjoint copies of $K_{12}$ in $G[U]$. By the cap $D(13)=67$, every vertex outside a target $K_{12}$ has at most one target neighbor in it: otherwise the $13$-set consisting of the clique and that vertex would span at least $66+2=68$ target edges. In particular, between any two packed blocks the target edges form a matching.

There cannot be $12$ packed blocks. Indeed, choose one representative from each block greedily. At the $k$th step, the previously chosen $k-1$ representatives exclude at most $k-1$ vertices of the current $K_{12}$, so a new representative can be chosen nonadjacent to all previous representatives. The twelve representatives, together with $v$, would form an independent $13$-set. Thus $0\leq j\leq11$.

Let $R$ be the induced graph on the vertices of $U$ outside the packed blocks. Maximality gives $\omega(R)\leq11$. Moreover,
\[
\alpha(R)\leq11-j.
\]
For if $R$ contained an independent set of size $12-j$, each packed block would lose at most $12-j$ vertices to adjacency with this set and would retain at least $j$ vertices. Greedily choose independent representatives from $j$ blocks and add $v$; this again gives an independent $13$-set. This argument includes $j=0$ and $j=11$.

\subsection{The exact edge ledger}

Set
\[
\sigma=870-e(G),\qquad
z=\binom d2-e(G[A]),\qquad
g=\sum_{a\in A}(d_G(a)-d).
\]
Let $x$ be the number of target edges between distinct packed blocks and let $y$ be the number of target edges from packed blocks to $R$. All five quantities are nonnegative. The induced closed neighborhood has
\[
e(G[N[v]])=\binom{d+1}{2}-z.
\]
Also,
\[
\sum_{a\in A}d_G(a)=d^2+g
=d+2e(G[A])+e_G(A,U),
\]
and hence
\[
e_G(A,U)=2z+g.
\]
Partitioning all edges according to $N[v]$, the packed cliques, and $R$ gives the exact identity
\begin{equation}
e(R)=870-\binom{d+1}{2}-66j-(\sigma+z+g+x+y).
\label{eq:frontend-budget}
\end{equation}
In particular, the initial neighborhood $A$ is not assumed to be a clique: its defect $z$ is explicitly charged. The five nonnegative terms in \eqref{eq:frontend-budget} are not interchangeable and none is omitted.

\subsection{The contextual R(4,49) obstruction}

The recurrence below is initially formulated in the contextual class of induced subgraphs of the fixed target graph $G$. We first prove that no induced member of $\mathcal R(4,49)$ can occur in the original full-coloring setting.

\begin{lemma}[No contextual $49$-vertex residual]
\label{lem:context49}
There is no $49$-vertex induced target-color subgraph $F$ in the assumed balanced coloring satisfying
\[
\alpha(F)\leq4,\quad \omega(F)\leq11,\quad
e(F[W])\leq D(|W|)\ \text{for every }W\subseteq V(F).
\]
\end{lemma}

\begin{proof}
First $\delta(F)\geq13$. A vertex of degree at most $11$ leaves at least $37$ nonneighbors, which contain an induced member of $\mathcal R(3,37)$, impossible by \cref{lem:R336}. If a vertex $p$ has degree $12$, its $36$-vertex nonneighborhood $H$ belongs to $\mathcal R(3,36)$. The lower-order theorem gives $e(H)\geq234$, while $D(36)=234$, so $e(H)=234$. The $12$-vertex shell cannot be a clique, since together with $p$ it would form a $K_{13}$. Choose a nonedge $ab$ in the shell. The union of the $H$-neighborhoods of $a$ and $b$ has at least $13$ vertices: otherwise at least $24$ vertices of $H$ avoid both, giving a member of $\mathcal R(2,24)$. Therefore $H\cup\{a,b\}$ spans at least $234+13=247>D(38)=241$ target edges, a contradiction. Thus $\delta(F)\geq13$ and $e(F)\geq\lceil49\cdot13/2\rceil=319$.

Now fix one of the other eleven colors, and let $C$ be its graph on the same $49$ vertices. Balance gives $\alpha(C)\leq12$. Also $\omega(C)\leq4$, since a $K_5$ in $C$ would be an independent five-set in the target graph $F$. The complement $\overline C$ is $K_{13}$-free but is not $12$-partite: if it were, its twelve independent parts would be twelve cliques in $C$, each of size at most four, and could cover at most $48$ vertices. Kang--Pikhurko's second-Tur\'an theorem applies with $n=49,r=12$ and gives
\[
e(\overline C)\leq t_{12}(49)-\left\lfloor\frac{49}{12}\right\rfloor+1
=1100-4+1=1097,
\]
where $t_{12}(49)=\ex(49,K_{13})=1100$. Hence $e(C)\geq1176-1097=79$. The eleven non-target colors therefore contribute at least $869$ edges on this $49$-set, so $e(F)\leq1176-869=307$, contradicting $e(F)\geq319$.
\end{proof}

The lemma is deliberately contextual. It does not prove that every abstract uncolored graph in $\mathcal R(4,49)$ is impossible; the eleven non-target color classes are part of the argument.

\subsection{The recurrence and its finite domain}

For an induced subgraph of the fixed target graph $G$ with independence number at most $a$ and clique number at most $11$, let $B(a,n)$ denote a proved lower bound on its number of edges. The base cases are $B(a,0)=0$, the class $a=0$ is empty at positive order, and $a=1$ consists of cliques of order at most $11$. The certified contextual order maxima are $M_0=0$, $M_1=11$, $M_2=23$, $M_3=36$, and $M_4=48$; for larger $a$, $M_a$ is the maximum order retained by the same recurrence table. Equivalently, the initial empty-class cutoffs are $24$ for $a=2$, $37$ for $a=3$, and $49$ for $a=4$. The seed lower bounds are
\[
B(2,23)\geq131,\quad B(3,35)\geq210,\quad
B(3,36)\geq234,\quad B(4,47)\geq284,\quad B(4,48)\geq308.
\]

For completeness, the recurrence rule is recorded here. In a graph $J$ of order $n$ and $\alpha(J)\leq a$, choose a minimum-degree vertex $p$ of degree $h$, put $A=N_J(p)$ and $C=J-N_J[p]$, and let $M=\overline{J[A]}$. For $x\in A$, set
\[
g_x=d_J(x)-h,\quad t_x=|N_C(x)|,\quad z=e(M),\quad
q=z+\sum_{x\in A}g_x.
\]
Then
\begin{align}
t_x&=d_M(x)+g_x,\\
e(J)&=\binom{h+1}{2}+q+e(C),\\
q&=t(S)-e_M(S)+e_M(A\setminus S)+g(A\setminus S)
\quad(S\subseteq A).
\label{eq:recurrence-pivot}
\end{align}
Since $\omega(J)\leq11$, $\alpha(M)\leq10$, and $\omega(M)=\alpha(J[A])\leq a$, the shell and closed-shell caps imply
\[
q\geq z\geq
\max\left\{\tau_{10}(h),\binom{h+1}{2}-D(h+1),
\binom h2-D(h)\right\}.
\]

For $h\geq11$, $M$ has an edge. If $M_b$ denotes the certified maximum order of a contextual graph with independence number at most $b$, the union of the two core rows corresponding to an edge has size at least
\[
\beta=\max\{0,m-M_{a-2}\},\qquad m=n-h-1.
\]
Thus $q\geq\beta-1$ and $e(C)+\beta\leq D(m+2)$. A maximum matching in $M$ has size at least
\[
k(h)=\max\left\{\left\lceil\frac{h-10}{2}\right\rceil,
1+\mathbf 1_{h\geq12}+\mathbf 1_{h\geq13}\right\}.
\]
The first term follows because the unmatched vertices of a maximum matching are independent and $\alpha(M)\leq10$. The additional small-order cases are as follows. At $h=11$, $M$ has an edge. At $h=12$, the cap gives $z\geq11$; a graph with matching number at most one is a star or a triangle with isolated vertices, and these alternatives have independence number at least $11$ or at most three edges, respectively. Thus $M$ has a two-matching. At $h=13$, the cap gives $z\geq22$, so \cref{lem:matching13} gives a three-matching. At $h=14$, the $15$-set cap gives $z\geq33$. Matching number at most one is impossible: such a graph is a star or a triangle plus isolated vertices, forcing independence number at least $12$ or at most three edges. Thus if the matching number were at most two, it would be exactly two; the four endpoints $S$ of a maximum two-matching and the remaining independent ten-set $I$ would have a bipartite cross graph with matching number at most two. By K\H{o}nig's theorem it has a vertex cover of size two. Depending on whether the cover contains two vertices of $S$, one from each side, or two vertices of $I$, the cross-edge count is at most $20$, $13$, or $8$, respectively. Adding at most six edges inside $S$ gives $z\leq26$, a contradiction. For $h\geq15$, the first term is already at least three. This proves the stated $k(h)$ in every case.

For $t$ matching edges, the $2t$ endpoints have total weight at least $t\beta$, while $\omega(M)\leq a$ bounds their internal missing edges by $\binom{2t}{2}-\tau_a(2t)$. Consequently
\[
q\geq t\beta-\binom{2t}{2}+\tau_a(2t).
\]
If $M$ is triangle-free, Mantel's theorem \cite{Mantel1907} improves this to $q\geq t\beta-t^2$. If $M$ contains a triangle, the three corresponding core rows have union at least
\[
\gamma=\max\{0,m-M_{a-3}\},
\]
so $q\geq\gamma-3$ and $e(C)+\gamma\leq D(m+3)$. For $a=2$ the triangle alternative is impossible. The recurrence takes the minimum of the surviving triangle-free and triangle-channel lower bounds.

For each viable minimum degree $h$, put
\[
L_h=\max\left\{\left\lceil\frac{nh}{2}\right\rceil,
B(a-1,m)+\binom{h+1}{2}+Q_h\right\},
\]
where $Q_h$ is the minimum of the licensed shell-channel floors just described. The node bound is the maximum of $\min_h L_h$, $\tau_a(n)$, any applicable seed, and the deletion-averaging bound
\[
\left\lceil\frac{nB(a,n-1)}{n-2}\right\rceil\quad(n>2).
\]
The last inequality follows by summing the edge counts over all one-vertex deletions, since each edge is counted exactly $n-2$ times. A node is empty if a predecessor is empty, no minimum-degree row survives, or the resulting lower bound exceeds $D(n)$. Every recursive reference has smaller $a$ or the same $a$ and smaller $n$, so the deduction is finite and acyclic.

Three independent integer implementations agree on all $1,752$ recurrence nodes and all $144$ outer-stop comparisons. The unique surviving stop is
\[
(d,j,a,n)=(11,6,5,61),\qquad B(5,61)=397,
\]
with edge budget $408$ from \eqref{eq:frontend-budget}. Thus a hypothetical coloring yields a single induced graph $F$ with
\begin{equation}
|V(F)|=61,\quad 397\leq E=e(F)\leq408,\quad
\delta(F)\in\{12,13\},\quad
\alpha(F)\leq5,\quad\omega(F)\leq11,
\label{eq:T61-contract}
\end{equation}
and every induced subset satisfies the cap $D$. Indeed, a vertex of degree at most $11$ would leave at least $49$ nonneighbors, any $49$ of which form a contextual member of $\mathcal R(4,49)$, contrary to \cref{lem:context49}. If $\delta(F)\geq14$, the handshake lemma gives $E\geq\lceil61\cdot14/2\rceil=427>408$.
